\begin{afterword}
\summaryhead

The essay distinguishes two kinds of nonconformity. Joining a rebellion that already exists is fairly easy. In Asch's conformity experiments a single ally made subjects far more willing to contradict the group. Being the first to dissent, alone, is much harder.

The author argues that ordinary rebellion, such as dressing in black or being vegetarian, is understood by the people around you. Real lonely dissent meets incomprehension instead of hostility, and that is more frightening. The example is cryonics. The fear of thinking really differently, which cryonics demands, is said to be stronger than the fear of death. The author offers an evolutionary explanation: being driven out of the band alone was once fatal.

The highest form is the scientific revolutionary, who reaches new territory by pursuing correct answers. The author adds that courage is no proof of being right, and that people who dissent easily, the author included, should correct for it. Nonconformity, the essay concludes, is not a virtue in itself.

\responsehead

The essay is presented as a lesson about courage. What it does is explain why other people do not share the author's unusual beliefs. The people who find cryonics strange are not given a reason. They are given noises (``Huh? What?''), a fear stronger than the fear of death, and an evolutionary history to account for the fear. The people who sign up are given courage. A reader who holds an unpopular belief is offered a flattering account of every blank look they receive, and the essay gives no way to tell that account from the plain one, that the belief looks wrong to others for reasons.

Each step has little support. The one experiment cited is reported selectively. In Asch's unanimous-majority condition, most answers were lonely dissent, and a quarter of subjects never conformed at all. The post calls its evolutionary story post facto and then uses it as a premise. The claim that fear of weirdness beats fear of death is based on comparing skydiving with a procedure that cost tens of thousands of dollars and has never been shown to work.

The essay divides its space unevenly. For most of its length the essay treats courage as the scarce thing that separates real dissent from fake. The two short paragraphs conceding that ``the difficulty'' is in being right, and that dissent is ``just a bias either way,'' are the ones a reader most needs, and they get the least space. The paragraph of self-description claims, with some irony, the rare trait the essay says most people only imagine they have; the next confirms that dissenting too easily is ``my own nature.'' The essay's own footnote points to Robin Hanson's argument that careless dissenters, not conformity, are the main obstacle to good new ideas. The last paragraph agrees with Hanson, but the essay does not apply the point to cryonics.

In short: a defense of the author's own unpopular cause presented as a lesson on courage, built on a selectively reported experiment and an unexamined comparison with skydiving, with its most important caveat given two short paragraphs.
\end{afterword}
